\clearpage\thispagestyle{empty}
\pdfbookmark[0]{Quellen zu Juno (31. Januar 1814)}{juno1814jan-sources}
\begin{center}\textbf{Quellen zu Juno (31. Januar 1814)}\end{center}
\noindent Diese KI-Übertragung ins heutige Deutsch folgt dem Bericht mit der Überschrift Juno in den \href{https://gdz.sub.uni-goettingen.de/id/PPN236007467}{\textit{Werken}, Band~VI} von Gauß (Göttingen, 1874), S.~375. Der Bericht erschien unter Göttingen in den \textit{Göttingischen gelehrten Anzeigen}, Stück~19 (31. Januar 1814), \href{https://api.digitale-sammlungen.de/iiif/image/v2/bsb10538685_00195/full/1600,/0/default.jpg}{S.~189}. Der Link führt zum digitalisierten Exemplar der Bayerischen Staatsbibliothek.
\par\smallskip\noindent In beiden Ausgaben steht über der Beobachtungstabelle das Jahr 1812. Die Beobachtungen in Bessels Fußnote im gesonderten Briefwechsel vom Dezember 1813 tragen die Jahresangabe 1813. KI-Übertragung und KI-Satz.
\clearpage
